Quantum machine learning (\qml) is widely promoted as a route to learning tasks
that classical machine learning cannot handle, yet the field's central
question---\emph{when does a quantum model actually beat the best classical
model?}---remains unsettled and frequently overstated. The honest answer is
two-sided: rigorous classical-hard/quantum-easy separations exist but require
fault-tolerant hardware and are not simulable at 20--30 qubits, while the only
demonstrable separations at simulable scale live on engineered tasks. We provide
a single, reproducible study that pins down both sides at \emph{20 qubits or
more} under exact statevector simulation. On an engineered task built along the
direction of maximal geometric difference between a quantum and a classical
kernel, a projected-quantum-kernel SVM reaches $0.865$ test accuracy at 20 qubits
while the best of six tuned classical models reaches $0.629$ (paired
$t$-test $p=3.9\times10^{-8}$, Cohen's $d=5.4$); the advantage is stable at
$\approx0.22$--$0.23$ across $8$, $10$, and $20$ qubits. A within-quantum control
shows the naive fidelity kernel collapses to chance ($0.493$) at 20 qubits,
confirming the win comes from the \emph{right} quantum inductive bias rather than
quantumness alone. On the discrete-logarithm task, six off-the-shelf classical
models stay at chance ($0.49$--$0.52$) for $12$/$17$/$21$-bit problems while the
quantum interval-state kernel separates the task near-perfectly
($0.997$--$1.000$). We document every caveat---engineered labels, asymptotic
hardness, and noiseless infinite-shot simulation---so the claim is confirmed with
honest scope: a 20+-qubit simulated quantum kernel generalizes on tasks where the
best classical models cannot.
